\section{Threats to Validity}
\label{sec:threats}

All measurements are on one GPU (\gpuModel{}); the fitted constants are
A100-specific but the shapes should transfer. The batching ceiling reproduces
across two runtimes and two model families, but both measurements are
engine-bounded (Appendix~\ref{app:batching}). The crossover is calibrated at
a $\sim$\refPrompt-token prompt and shifts with input length. The AR--DG
quality difference ($\Delta = \qualityGap$) has a 95\% CI of $[-0.012,
+0.011]$ at $n=300$; the two are statistically equivalent at a $\pm 0.01$
margin.

The surviving niche has not been tested on a task whose natural output length
is $\geq 512$ tokens. The forced comparison in Table~\ref{tab:niche} charges
the AR arm for tokens that first-object-wins parsing discards, so the
$5.4\times$ advantage is an upper bound on AR cost rather than a measurement
of it. A task with genuinely long natural output --- multi-record extraction,
incident reports --- would test the niche without this artifact.

The niche has also not been tested against grammar-constrained AR decoding
(outlines, xgrammar, llguidance), which would eliminate the over-generation
the forced comparison charges for, guarantee 100\% schema validity, and
remove the diffusion model's completeness advantage. Constrained decoding is
the standard answer to the structured-output problem and is a stated
limitation of this paper.
